% Formato de Artículo de "Almanaque Económico" en Español

\documentclass[utf8, twoside,]{article}

%Paquetes bàsicos-----------------------
\usepackage[utf8]{inputenc}
\usepackage{standalone}
\usepackage{pgfplots}
    \pgfplotsset{compat=1.15}
\usepackage{tikz}
\usepackage{import}
\usepackage{amsmath}
\usepackage[autostyle=true]{csquotes}
\usepackage[american,spanish]{babel}
\usepackage{biblatex-chicago}
\usepackage{csquotes}
\usepackage{graphicx}
\usepackage{setspace} 
%---------------------------------------

% Directorio de imágenes. Las imágenes deben ir en la carpeta "figs". Las tablas deben ir en la carpeta "tables".

\graphicspath{{figs/}}

%---------------------------------------

% Se crean comandos de palabras clave---
\providecommand{\keywords}[1]
{
  \small	
  \textbf{\textit{Keywords---}} #1
}

\providecommand{\pclave}[1]
{
  \small	
  \textbf{\textit{Palabras clave---}} #1
}

% Establece cómo escribir "editado por" en las citas en español 

\NewBibliographyString{cbyeditor}
\DefineBibliographyStrings{spanish}{cbyeditor = {editado por}}

%Espaciado doble

\doublespacing

%---------------------------------------

% NO MODIFICAR NADA ANTES DE ESTE PUNTO

%Añade paquetes adicionales aquí--------



%---------------------------------------


%Añade tu bibliografía como archivo .bib 

\addbibresource{mybib.bib} 


%--------------DOCUMENTO-----------------


\begin{document} %Comienza el documento
\begin{otherlanguage}{spanish}

%Se añaden los datos del título del artículo

\title{Title1\\ \medskip %titulo en otro idioma en otra línea
Título1}
\author{Nombre Apellido, \it{Institución Educativa}}
\date{}

%Espaciado sencillo para primera página

\begin{singlespace}



\maketitle %Se crea el tìtulo



%Se define el nombre de la sección Abstract en inglés

\def\abstractname{{Abstract}}

%Se crea el abstract en inglés

\begin{abstract}
     \vspace{-0.2cm} %Ajuste de espacio en resumen
     
Texto del resumen en inglés
  
\end{abstract}

%Se define el nombre de la sección Abstract en español

\def\abstractname{{Resumen}}

%Se crea el abstract en español

\begin{abstract}
      \vspace{-0.2cm} %Ajuste de espacio en resumen
      
Texto del resumen en español
   
\end{abstract}

%Se introducen las palabras clave en inglés

\keywords{one, two, three}

%Se introducen las palabras clave en español

\pclave{uno, dos, tres}

%Cerramos el formato de página de Título



\end{singlespace}

\pagebreak

%Introducción

\section{Introducción}

Texto de la introducción. Ejemplo de uso de citas:


Lorem ipsum dolor sit amet, consectetur adipiscing elit.\footnote{\cite{fernandez}} Pellentesque tincidunt sit amet magna a ultricies. Praesent ultrices dictum velit, ut luctus sem fermentum a.\footnote{\cite{anton}} Aenean sed ex fringilla nulla mattis cursus. Sed efficitur ligula condimentum purus vestibulum placerat.\footnote{\cite{fernandez}}



%Cuerpo

\section{Cuerpo}
  
Texto del cuerpo. Se puede renombrar el título de esta sección y crear secciones adicionales.

% Ejemplo de imagen. El formate debe ser .eps
% La imágen de ejemplo fue tomada de https://people.sc.fsu.edu/~jburkardt/data/eps/eps.html

\begin{figure}[ht]
\includegraphics[]{mathematica.eps}
\caption{Imagen eps} \label{fig:img1}
\end{figure}

%Aceptamos gráficas en formato .eps, pero es preferible crear la gráfica en Latex usando el paquete pgfplots o TikZ para mayor consistencia

% Imagen con pgfplots

Agregamos una imágen creada con pgfplots:


\begin{figure}[ht]
\subimport{figs/}{graph.tex}
\caption{Imagen pgfplots} \label{fig:img2}
\end{figure}

%Imágen con TikZ

Ejemplo de gráfica usando TikZ

\begin{figure}[ht]
\subimport{figs/}{graph2.tex}
\caption{Imagen TikZ} \label{fig:img3}
\end{figure}

Usa el método más apto para tus necesidades.

También es posible hacer referencia a una imagen o tabla. La Figura \ref{fig:img1}, por ejemplo, es una imagen en formato .eps





 %Agregamos una tabla de regresión. Para ver cómo exportar una tabla de regresión de Stata a Latex, consultar la guía en la página web
 % http://repec.org/bocode/e/estout/esttab.html
 
 Una tabla de regresión
 
 \begin{table}[ht]
 \subimport{tables/}{reg.tex}
 \caption{Tabla importada de Stata} \label{table:reg1}
\end{table}

%Ejemplo de matemáticas
 
 Ejemplo de matemáticas
 
 La matriz de sustitución de Slutsky de 2 bienes es:
 
 \begin{equation}\label{eq:1}
    \begin{bmatrix}
 \dfrac{\delta^2e(p.u)}{\delta p_1^2} & \dfrac{\delta^2e(p.u)}{\delta p_2 \delta p_1} \\[2ex]
 \dfrac{\delta^2e(p.u)}{\delta p_1\delta p_2} & \dfrac{\delta^2e(p.u)}{\delta p_2^2} \\[2ex]
 \end{bmatrix}
\end{equation}
 

  Y por el Teorema de Shepard sabemos que 
  
   \begin{equation}\label{eq:2}
   \dfrac{\delta e (\vec{p}, u)}{\delta p_i} = x_i^h(\vec{p}, u)
   \end{equation}
   
   Donde $x_i^h(\vec{p}, u))$ es la demanda Hicksiana,
  
   Ahora, por \eqref{eq:2}, podemos escribir \eqref{eq:1} como:
   
    \begin{equation}
    \begin{bmatrix}
 \dfrac{\delta x_1^h(\vec{p}, u)}{\delta p_1} 
 & \dfrac{\delta x_1^h(\vec{p}, u)}{\delta p_2} \\[2ex]
 \dfrac{\delta x_2^h(\vec{p}, u)}{\delta p_1} 
 & \dfrac{\delta x_2^h(\vec{p}, u)}{\delta p_2} \\[2ex]
 \end{bmatrix}
\end{equation}
  
   La tercera ley de Hicks es:
   \begin{equation*}
      \sum_{j=1}^n \dfrac{\delta x_i^h(\vec{p}, u)}{\delta p_j} p_j = 0, \quad i=1,\dots,n
  \end{equation*}
  
  Por lo que
  
      \begin{equation*}
    \begin{bmatrix}
 \dfrac{\delta x_1^h(\vec{p}, u)}{\delta p_1} 
 & \dfrac{\delta x_1^h(\vec{p}, u)}{\delta p_2} \\[2ex]
 \dfrac{\delta x_2^h(\vec{p}, u)}{\delta p_1} 
 & \dfrac{\delta x_2^h(\vec{p}, u)}{\delta p_2} \\[2ex]
 \end{bmatrix}
 \begin{bmatrix}
 p_1 \\[2ex]
 p_2 \\[2ex]
 \end{bmatrix}
 =
 \begin{bmatrix}
 0 \\[2ex]
 0 \\[2ex]
 \end{bmatrix}
\end{equation*} 
 
%Conclusiones

\section{Conclusiones}

Texto de las conclusiones.

\pagebreak

%Se define el idioma de la bibliografía


\printbibliography %inserta la bibliografía
\vfill

\end{otherlanguage}


\end{document}